% !TeX program = lualatex
% ============================================================
%  brief-layout.tex – gemeinsames Layout für alle Briefe
%  Formaler deutscher Brief nach DIN 5008 (Form B)
%  Klasse: KOMA-Script scrlttr2 · Engine: LuaLaTeX
%
%  Liegt im Ordner _vorlage/ und wird von JEDEM Brief in JEDEM
%  Themenordner mit % !TeX program = lualatex
% ============================================================
%  brief-layout.tex – gemeinsames Layout für alle Briefe
%  Formaler deutscher Brief nach DIN 5008 (Form B)
%  Klasse: KOMA-Script scrlttr2 · Engine: LuaLaTeX
%
%  Liegt im Ordner _vorlage/ und wird von JEDEM Brief in JEDEM
%  Themenordner mit % !TeX program = lualatex
% ============================================================
%  brief-layout.tex – gemeinsames Layout für alle Briefe
%  Formaler deutscher Brief nach DIN 5008 (Form B)
%  Klasse: KOMA-Script scrlttr2 · Engine: LuaLaTeX
%
%  Liegt im Ordner _vorlage/ und wird von JEDEM Brief in JEDEM
%  Themenordner mit % !TeX program = lualatex
% ============================================================
%  brief-layout.tex – gemeinsames Layout für alle Briefe
%  Formaler deutscher Brief nach DIN 5008 (Form B)
%  Klasse: KOMA-Script scrlttr2 · Engine: LuaLaTeX
%
%  Liegt im Ordner _vorlage/ und wird von JEDEM Brief in JEDEM
%  Themenordner mit \input{brief-layout} geladen. Gefunden wird es über
%  _vorlage/latexmkrc (setzt die Suchpfade für TeX-Dateien und Schriften).
%  Schriften: _vorlage/fonts/ · Absenderdaten: <Person>/absender.lco
% ============================================================
\documentclass[
  fontsize=11pt,
  paper=a4,
  parskip=half,
  enlargefirstpage,
  foldmarks=true,
  foldmarks=TBMP,        % Faltmarken oben/Mitte/unten + Lochmarke
  backaddress=true,      % Rücksendeangabe im Fenster
  subject=beforeopening,
  subject=untitled,      % Betreff fett, ohne "Betreff:" (DIN 5008)
]{scrlttr2}

\LoadLetterOption{DIN}   % Anschriftfeld/Geometrie für DIN-lang-Fensterkuvert

% --- Satzspiegel nach DIN 5008: links 25 mm, rechts 20 mm ------------------
% (KOMA-Standard ergab ~31,5 mm links/rechts; Anschrift und Text standen
%  dadurch nicht auf derselben Fluchtlinie.)
\setlength{\oddsidemargin}{\dimexpr 25mm-1in\relax}
\setlength{\evensidemargin}{\oddsidemargin}
\setlength{\textwidth}{\dimexpr\paperwidth-25mm-20mm\relax}
% Anschrift beginnt bei 20 mm (Fensterrand) + 5 mm Einzug = 25 mm,
% also bündig mit Brieftext, Kopf und Betreff.
\setplength{toaddrindent}{5mm}
% Unterer Rand der Folgeseiten: KOMA-Standard ließ den Text bei ca. 239 mm
% enden (fast 6 cm Rand). Jetzt endet er bei 265 mm, die Fußzeile steht bei
% 277 mm. (Seite 1 wird über enlargefirstpage bis kurz vor den Fuß genutzt.)
\setlength{\textheight}{\dimexpr 265mm-(1in+\voffset+\topmargin+\headheight+\headsep)\relax}

% --- Grafiken (eingescannte Unterschrift) — vor absender.lco laden! ------
\usepackage{graphicx}

% --- Absenderdaten: nächstgelegene absender.lco -----------------------------
% Gesucht wird vom Ordner des Briefes aus nach oben. So gilt Max/absender.lco
% für alles unter Max/, Erika/absender.lco für alles unter Erika/.
% Ein Brief kann eine eigene absender.lco im eigenen Ordner haben (gewinnt).
% Nur ein Absender? Dann genügt eine absender.lco im Projektordner.
\directlua{
  local d = lfs.currentdir():gsub(string.char(92), "/")  % Windows: \ -> /
  local found = ""
  while d and d ~= "" do
    if lfs.isfile(d .. "/absender.lco") then found = d break end
    local parent = d:match("^(.*)/[^/]+$")
    if not parent or parent == "" or parent == d then break end
    d = parent
  end
  token.set_macro("absenderdir", found)
}
\ifx\absenderdir\empty
  \errmessage{Keine absender.lco gefunden: Lege eine absender.lco in den
    Ordner des Briefes oder einen Ordner darüber (Vorlage:
    _vorlage/absender-vorlage.lco)}
\else
  \input{\absenderdir/absender.lco}
\fi

% --- Schrift: Source Sans 3 aus ./fonts/ ---------------------------------
\usepackage{fontspec}
% Schriften werden per Dateiname gefunden (Suchpfad aus latexmkrc),
% darum hier KEIN Path=fonts/ – so klappt es aus jedem Unterordner.
\setmainfont{SourceSans3}[
  Extension      = .otf,
  UprightFont    = *-Regular,
  ItalicFont     = *-It,
  BoldFont       = *-Semibold,       % Semibold als \textbf: ruhiger als Bold
  BoldItalicFont = *-SemiboldIt,
  Scale = 1.08,
]
\setsansfont{SourceSans3}[
  Extension=.otf,
  UprightFont=*-Regular, ItalicFont=*-It,
  BoldFont=*-Semibold, BoldItalicFont=*-SemiboldIt,
]% auch \sffamily (z. B. Rücksendeangabe) nutzt Source Sans
\newfontfamily\NameSchrift{PlayfairDisplaySC-Black}[Extension=.ttf]

\linespread{1.06}

\usepackage[ngerman]{babel}
\usepackage{microtype}

% Kompakte Aufzählungen (z. B. Fragenkataloge)
\usepackage{enumitem}
\setlist{itemsep=0pt, parsep=0pt, topsep=2pt, partopsep=0pt, leftmargin=1.6em}

% --- Farben ---------------------------------------------------------------
\usepackage{xcolor}
\definecolor{akzent}{RGB}{0, 82, 100}    % dunkles Petrol
\definecolor{grau}{RGB}{88, 92, 96}      % Text im Kopf/Fuß

% --- PDF-Metadaten + anklickbare Links (ohne farbige Rahmen) ---------------
% Titel setzt jeder Brief selbst: \hypersetup{pdftitle={...}}
\usepackage[hidelinks]{hyperref}
\hypersetup{pdfauthor={\absendername}, pdflang={de-DE}}

% --- Seitenzahl "1/2" — nur bei mehrseitigen Briefen (DIN 5008) ------------
% scrlttr2 schreibt für jeden Brief das Label "<Briefnr>.lastpage".
% \thisletter liefert die Nummer des aktuellen Briefes, damit funktioniert
% die Zählung auch, wenn mehrere Briefe in einer Datei stehen.
% refcount liest die Seitenzahl als Zahl aus; ist der Brief einseitig,
% bleibt die Fußzeile leer.
\usepackage{refcount}
\newcommand*{\seitenzahl}{%
  \ifnum\getpagerefnumber{\thisletter.lastpage}>1
    \thepage/\pageref{\thisletter.lastpage}%
  \fi}

% Folgeseiten: Seitenzahl unten LINKS (KOMA-Mechanismus scrlayer-scrpage)
\usepackage[automark]{scrlayer-scrpage}
\clearpairofpagestyles
\ifoot*{\footnotesize\color{grau}\seitenzahl}
\pagestyle{scrheadings}

% ==========================================================================
%  BRIEFKOPF: Name links, Kontaktblock rechts, feine Akzentlinie darunter.
%  Gespeist aus den Variablen in absender.lco – hier nichts eintragen.
% ==========================================================================
% Kopf/Fuß der ersten Seite exakt auf den Satzspiegel ausrichten. Erst bei
% \begin{document} setzen, weil KOMA diese Längen vorher neu berechnet.
\AtBeginDocument{%
  \setplength{firstheadwidth}{\textwidth}%
  \setplength{firstheadhpos}{\dimexpr 1in+\hoffset+\oddsidemargin\relax}%
  \setplength{firstfootwidth}{\textwidth}%
  \setplength{firstfoothpos}{\dimexpr 1in+\hoffset+\oddsidemargin\relax}%
}
\setplength{firstfootvpos}{277mm}

\setkomavar{firsthead}{%
  \begin{minipage}[b]{0.52\textwidth}
    {\NameSchrift\fontsize{18pt}{20pt}\selectfont\color{akzent}%
      \usekomavar{fromname}}%
  \end{minipage}%
  \hfill
  \begin{minipage}[b]{0.44\textwidth}
    \raggedleft\footnotesize\color{grau}%
    \usekomavar{fromaddress}\\[1pt]
    \usekomavar{fromphone}\\[1pt]
    \usekomavar{fromemail}%
  \end{minipage}\\[2.5mm]
  {\color{akzent}\rule{\textwidth}{0.6pt}}%
}

% Erste Seite: Seitenzahl ebenfalls unten links (nur wenn mehrseitig)
\setkomavar{firstfoot}{\footnotesize\color{grau}\seitenzahl\hfill}

% --- Bezugszeichen: nur setzen, wenn es eins gibt (pro Brief) --------------
% z. B. \setkomavar{yourref}{AB-12345}     -> "Ihr Zeichen"
%       \setkomavar{customer}[Aktenzeichen:~]{12345}

% --- Grußformel und Unterschrift linksbündig am Satzspiegel ----------------
\setplength{sigindent}{0mm}
\renewcommand*{\raggedsignature}{\raggedright}
% Abstand Gruß -> Unterschrift. Der Platz für die Unterschrift selbst
% (Bild oder 15 mm Leerraum) kommt aus \unterschriftsbild in absender.lco.
\setplength{sigbeforevskip}{2mm}

\setkomavar{date}{\today}

% --- Anlagen: PDF oder Bild (JPG/PNG) hinten an den Brief anhängen ---------
% Nach \end{letter}, vor \end{document}:
%   \anlage[Mietvertrag]{anlagen/mietvertrag.pdf}   alle Seiten eines PDFs
%   \anlage[Ausweis]{anlagen/ausweis.jpg}            Foto oder Scan
%   \anlage[Rechnung][pages={1,3}]{anlagen/re.pdf}   nur bestimmte Seiten
% Jede Seite bekommt oben rechts den Stempel „Anlage N: <Titel>“.
% Die Anlagenliste im Brief selbst bleibt \encl{...} vor \end{letter}.
\usepackage{pdfpages}
\newcounter{anlage}
\NewDocumentCommand{\anlage}{o O{} m}{%
  \stepcounter{anlage}%
  \def\anlagestempel{Anlage~\theanlage\IfValueT{#1}{: #1}}%
  \IfFileExists{#3}{%
    \includepdf[pages=-, scale=0.92, #2,
      pagecommand={\thispagestyle{empty}},
      picturecommand={\put(\LenToUnit{\paperwidth-20mm},\LenToUnit{\paperheight-12mm}){%
        \makebox[0pt][r]{\footnotesize\sffamily\color{grau}\anlagestempel}}}%
    ]{#3}%
  }{%
    \clearpage\thispagestyle{empty}%
    {\color{red}\bfseries \anlagestempel\ – Datei „\detokenize{#3}“ nicht gefunden}%
    \clearpage
  }%
}
 geladen. Gefunden wird es über
%  _vorlage/latexmkrc (setzt die Suchpfade für TeX-Dateien und Schriften).
%  Schriften: _vorlage/fonts/ · Absenderdaten: <Person>/absender.lco
% ============================================================
\documentclass[
  fontsize=11pt,
  paper=a4,
  parskip=half,
  enlargefirstpage,
  foldmarks=true,
  foldmarks=TBMP,        % Faltmarken oben/Mitte/unten + Lochmarke
  backaddress=true,      % Rücksendeangabe im Fenster
  subject=beforeopening,
  subject=untitled,      % Betreff fett, ohne "Betreff:" (DIN 5008)
]{scrlttr2}

\LoadLetterOption{DIN}   % Anschriftfeld/Geometrie für DIN-lang-Fensterkuvert

% --- Satzspiegel nach DIN 5008: links 25 mm, rechts 20 mm ------------------
% (KOMA-Standard ergab ~31,5 mm links/rechts; Anschrift und Text standen
%  dadurch nicht auf derselben Fluchtlinie.)
\setlength{\oddsidemargin}{\dimexpr 25mm-1in\relax}
\setlength{\evensidemargin}{\oddsidemargin}
\setlength{\textwidth}{\dimexpr\paperwidth-25mm-20mm\relax}
% Anschrift beginnt bei 20 mm (Fensterrand) + 5 mm Einzug = 25 mm,
% also bündig mit Brieftext, Kopf und Betreff.
\setplength{toaddrindent}{5mm}
% Unterer Rand der Folgeseiten: KOMA-Standard ließ den Text bei ca. 239 mm
% enden (fast 6 cm Rand). Jetzt endet er bei 265 mm, die Fußzeile steht bei
% 277 mm. (Seite 1 wird über enlargefirstpage bis kurz vor den Fuß genutzt.)
\setlength{\textheight}{\dimexpr 265mm-(1in+\voffset+\topmargin+\headheight+\headsep)\relax}

% --- Grafiken (eingescannte Unterschrift) — vor absender.lco laden! ------
\usepackage{graphicx}

% --- Absenderdaten: nächstgelegene absender.lco -----------------------------
% Gesucht wird vom Ordner des Briefes aus nach oben. So gilt Max/absender.lco
% für alles unter Max/, Erika/absender.lco für alles unter Erika/.
% Ein Brief kann eine eigene absender.lco im eigenen Ordner haben (gewinnt).
% Nur ein Absender? Dann genügt eine absender.lco im Projektordner.
\directlua{
  local d = lfs.currentdir():gsub(string.char(92), "/")  % Windows: \ -> /
  local found = ""
  while d and d ~= "" do
    if lfs.isfile(d .. "/absender.lco") then found = d break end
    local parent = d:match("^(.*)/[^/]+$")
    if not parent or parent == "" or parent == d then break end
    d = parent
  end
  token.set_macro("absenderdir", found)
}
\ifx\absenderdir\empty
  \errmessage{Keine absender.lco gefunden: Lege eine absender.lco in den
    Ordner des Briefes oder einen Ordner darüber (Vorlage:
    _vorlage/absender-vorlage.lco)}
\else
  \input{\absenderdir/absender.lco}
\fi

% --- Schrift: Source Sans 3 aus ./fonts/ ---------------------------------
\usepackage{fontspec}
% Schriften werden per Dateiname gefunden (Suchpfad aus latexmkrc),
% darum hier KEIN Path=fonts/ – so klappt es aus jedem Unterordner.
\setmainfont{SourceSans3}[
  Extension      = .otf,
  UprightFont    = *-Regular,
  ItalicFont     = *-It,
  BoldFont       = *-Semibold,       % Semibold als \textbf: ruhiger als Bold
  BoldItalicFont = *-SemiboldIt,
  Scale = 1.08,
]
\setsansfont{SourceSans3}[
  Extension=.otf,
  UprightFont=*-Regular, ItalicFont=*-It,
  BoldFont=*-Semibold, BoldItalicFont=*-SemiboldIt,
]% auch \sffamily (z. B. Rücksendeangabe) nutzt Source Sans
\newfontfamily\NameSchrift{PlayfairDisplaySC-Black}[Extension=.ttf]

\linespread{1.06}

\usepackage[ngerman]{babel}
\usepackage{microtype}

% Kompakte Aufzählungen (z. B. Fragenkataloge)
\usepackage{enumitem}
\setlist{itemsep=0pt, parsep=0pt, topsep=2pt, partopsep=0pt, leftmargin=1.6em}

% --- Farben ---------------------------------------------------------------
\usepackage{xcolor}
\definecolor{akzent}{RGB}{0, 82, 100}    % dunkles Petrol
\definecolor{grau}{RGB}{88, 92, 96}      % Text im Kopf/Fuß

% --- PDF-Metadaten + anklickbare Links (ohne farbige Rahmen) ---------------
% Titel setzt jeder Brief selbst: \hypersetup{pdftitle={...}}
\usepackage[hidelinks]{hyperref}
\hypersetup{pdfauthor={\absendername}, pdflang={de-DE}}

% --- Seitenzahl "1/2" — nur bei mehrseitigen Briefen (DIN 5008) ------------
% scrlttr2 schreibt für jeden Brief das Label "<Briefnr>.lastpage".
% \thisletter liefert die Nummer des aktuellen Briefes, damit funktioniert
% die Zählung auch, wenn mehrere Briefe in einer Datei stehen.
% refcount liest die Seitenzahl als Zahl aus; ist der Brief einseitig,
% bleibt die Fußzeile leer.
\usepackage{refcount}
\newcommand*{\seitenzahl}{%
  \ifnum\getpagerefnumber{\thisletter.lastpage}>1
    \thepage/\pageref{\thisletter.lastpage}%
  \fi}

% Folgeseiten: Seitenzahl unten LINKS (KOMA-Mechanismus scrlayer-scrpage)
\usepackage[automark]{scrlayer-scrpage}
\clearpairofpagestyles
\ifoot*{\footnotesize\color{grau}\seitenzahl}
\pagestyle{scrheadings}

% ==========================================================================
%  BRIEFKOPF: Name links, Kontaktblock rechts, feine Akzentlinie darunter.
%  Gespeist aus den Variablen in absender.lco – hier nichts eintragen.
% ==========================================================================
% Kopf/Fuß der ersten Seite exakt auf den Satzspiegel ausrichten. Erst bei
% \begin{document} setzen, weil KOMA diese Längen vorher neu berechnet.
\AtBeginDocument{%
  \setplength{firstheadwidth}{\textwidth}%
  \setplength{firstheadhpos}{\dimexpr 1in+\hoffset+\oddsidemargin\relax}%
  \setplength{firstfootwidth}{\textwidth}%
  \setplength{firstfoothpos}{\dimexpr 1in+\hoffset+\oddsidemargin\relax}%
}
\setplength{firstfootvpos}{277mm}

\setkomavar{firsthead}{%
  \begin{minipage}[b]{0.52\textwidth}
    {\NameSchrift\fontsize{18pt}{20pt}\selectfont\color{akzent}%
      \usekomavar{fromname}}%
  \end{minipage}%
  \hfill
  \begin{minipage}[b]{0.44\textwidth}
    \raggedleft\footnotesize\color{grau}%
    \usekomavar{fromaddress}\\[1pt]
    \usekomavar{fromphone}\\[1pt]
    \usekomavar{fromemail}%
  \end{minipage}\\[2.5mm]
  {\color{akzent}\rule{\textwidth}{0.6pt}}%
}

% Erste Seite: Seitenzahl ebenfalls unten links (nur wenn mehrseitig)
\setkomavar{firstfoot}{\footnotesize\color{grau}\seitenzahl\hfill}

% --- Bezugszeichen: nur setzen, wenn es eins gibt (pro Brief) --------------
% z. B. \setkomavar{yourref}{AB-12345}     -> "Ihr Zeichen"
%       \setkomavar{customer}[Aktenzeichen:~]{12345}

% --- Grußformel und Unterschrift linksbündig am Satzspiegel ----------------
\setplength{sigindent}{0mm}
\renewcommand*{\raggedsignature}{\raggedright}
% Abstand Gruß -> Unterschrift. Der Platz für die Unterschrift selbst
% (Bild oder 15 mm Leerraum) kommt aus \unterschriftsbild in absender.lco.
\setplength{sigbeforevskip}{2mm}

\setkomavar{date}{\today}

% --- Anlagen: PDF oder Bild (JPG/PNG) hinten an den Brief anhängen ---------
% Nach \end{letter}, vor \end{document}:
%   \anlage[Mietvertrag]{anlagen/mietvertrag.pdf}   alle Seiten eines PDFs
%   \anlage[Ausweis]{anlagen/ausweis.jpg}            Foto oder Scan
%   \anlage[Rechnung][pages={1,3}]{anlagen/re.pdf}   nur bestimmte Seiten
% Jede Seite bekommt oben rechts den Stempel „Anlage N: <Titel>“.
% Die Anlagenliste im Brief selbst bleibt \encl{...} vor \end{letter}.
\usepackage{pdfpages}
\newcounter{anlage}
\NewDocumentCommand{\anlage}{o O{} m}{%
  \stepcounter{anlage}%
  \def\anlagestempel{Anlage~\theanlage\IfValueT{#1}{: #1}}%
  \IfFileExists{#3}{%
    \includepdf[pages=-, scale=0.92, #2,
      pagecommand={\thispagestyle{empty}},
      picturecommand={\put(\LenToUnit{\paperwidth-20mm},\LenToUnit{\paperheight-12mm}){%
        \makebox[0pt][r]{\footnotesize\sffamily\color{grau}\anlagestempel}}}%
    ]{#3}%
  }{%
    \clearpage\thispagestyle{empty}%
    {\color{red}\bfseries \anlagestempel\ – Datei „\detokenize{#3}“ nicht gefunden}%
    \clearpage
  }%
}
 geladen. Gefunden wird es über
%  _vorlage/latexmkrc (setzt die Suchpfade für TeX-Dateien und Schriften).
%  Schriften: _vorlage/fonts/ · Absenderdaten: <Person>/absender.lco
% ============================================================
\documentclass[
  fontsize=11pt,
  paper=a4,
  parskip=half,
  enlargefirstpage,
  foldmarks=true,
  foldmarks=TBMP,        % Faltmarken oben/Mitte/unten + Lochmarke
  backaddress=true,      % Rücksendeangabe im Fenster
  subject=beforeopening,
  subject=untitled,      % Betreff fett, ohne "Betreff:" (DIN 5008)
]{scrlttr2}

\LoadLetterOption{DIN}   % Anschriftfeld/Geometrie für DIN-lang-Fensterkuvert

% --- Satzspiegel nach DIN 5008: links 25 mm, rechts 20 mm ------------------
% (KOMA-Standard ergab ~31,5 mm links/rechts; Anschrift und Text standen
%  dadurch nicht auf derselben Fluchtlinie.)
\setlength{\oddsidemargin}{\dimexpr 25mm-1in\relax}
\setlength{\evensidemargin}{\oddsidemargin}
\setlength{\textwidth}{\dimexpr\paperwidth-25mm-20mm\relax}
% Anschrift beginnt bei 20 mm (Fensterrand) + 5 mm Einzug = 25 mm,
% also bündig mit Brieftext, Kopf und Betreff.
\setplength{toaddrindent}{5mm}
% Unterer Rand der Folgeseiten: KOMA-Standard ließ den Text bei ca. 239 mm
% enden (fast 6 cm Rand). Jetzt endet er bei 265 mm, die Fußzeile steht bei
% 277 mm. (Seite 1 wird über enlargefirstpage bis kurz vor den Fuß genutzt.)
\setlength{\textheight}{\dimexpr 265mm-(1in+\voffset+\topmargin+\headheight+\headsep)\relax}

% --- Grafiken (eingescannte Unterschrift) — vor absender.lco laden! ------
\usepackage{graphicx}

% --- Absenderdaten: nächstgelegene absender.lco -----------------------------
% Gesucht wird vom Ordner des Briefes aus nach oben. So gilt Max/absender.lco
% für alles unter Max/, Erika/absender.lco für alles unter Erika/.
% Ein Brief kann eine eigene absender.lco im eigenen Ordner haben (gewinnt).
% Nur ein Absender? Dann genügt eine absender.lco im Projektordner.
\directlua{
  local d = lfs.currentdir():gsub(string.char(92), "/")  % Windows: \ -> /
  local found = ""
  while d and d ~= "" do
    if lfs.isfile(d .. "/absender.lco") then found = d break end
    local parent = d:match("^(.*)/[^/]+$")
    if not parent or parent == "" or parent == d then break end
    d = parent
  end
  token.set_macro("absenderdir", found)
}
\ifx\absenderdir\empty
  \errmessage{Keine absender.lco gefunden: Lege eine absender.lco in den
    Ordner des Briefes oder einen Ordner darüber (Vorlage:
    _vorlage/absender-vorlage.lco)}
\else
  \input{\absenderdir/absender.lco}
\fi

% --- Schrift: Source Sans 3 aus ./fonts/ ---------------------------------
\usepackage{fontspec}
% Schriften werden per Dateiname gefunden (Suchpfad aus latexmkrc),
% darum hier KEIN Path=fonts/ – so klappt es aus jedem Unterordner.
\setmainfont{SourceSans3}[
  Extension      = .otf,
  UprightFont    = *-Regular,
  ItalicFont     = *-It,
  BoldFont       = *-Semibold,       % Semibold als \textbf: ruhiger als Bold
  BoldItalicFont = *-SemiboldIt,
  Scale = 1.08,
]
\setsansfont{SourceSans3}[
  Extension=.otf,
  UprightFont=*-Regular, ItalicFont=*-It,
  BoldFont=*-Semibold, BoldItalicFont=*-SemiboldIt,
]% auch \sffamily (z. B. Rücksendeangabe) nutzt Source Sans
\newfontfamily\NameSchrift{PlayfairDisplaySC-Black}[Extension=.ttf]

\linespread{1.06}

\usepackage[ngerman]{babel}
\usepackage{microtype}

% Kompakte Aufzählungen (z. B. Fragenkataloge)
\usepackage{enumitem}
\setlist{itemsep=0pt, parsep=0pt, topsep=2pt, partopsep=0pt, leftmargin=1.6em}

% --- Farben ---------------------------------------------------------------
\usepackage{xcolor}
\definecolor{akzent}{RGB}{0, 82, 100}    % dunkles Petrol
\definecolor{grau}{RGB}{88, 92, 96}      % Text im Kopf/Fuß

% --- PDF-Metadaten + anklickbare Links (ohne farbige Rahmen) ---------------
% Titel setzt jeder Brief selbst: \hypersetup{pdftitle={...}}
\usepackage[hidelinks]{hyperref}
\hypersetup{pdfauthor={\absendername}, pdflang={de-DE}}

% --- Seitenzahl "1/2" — nur bei mehrseitigen Briefen (DIN 5008) ------------
% scrlttr2 schreibt für jeden Brief das Label "<Briefnr>.lastpage".
% \thisletter liefert die Nummer des aktuellen Briefes, damit funktioniert
% die Zählung auch, wenn mehrere Briefe in einer Datei stehen.
% refcount liest die Seitenzahl als Zahl aus; ist der Brief einseitig,
% bleibt die Fußzeile leer.
\usepackage{refcount}
\newcommand*{\seitenzahl}{%
  \ifnum\getpagerefnumber{\thisletter.lastpage}>1
    \thepage/\pageref{\thisletter.lastpage}%
  \fi}

% Folgeseiten: Seitenzahl unten LINKS (KOMA-Mechanismus scrlayer-scrpage)
\usepackage[automark]{scrlayer-scrpage}
\clearpairofpagestyles
\ifoot*{\footnotesize\color{grau}\seitenzahl}
\pagestyle{scrheadings}

% ==========================================================================
%  BRIEFKOPF: Name links, Kontaktblock rechts, feine Akzentlinie darunter.
%  Gespeist aus den Variablen in absender.lco – hier nichts eintragen.
% ==========================================================================
% Kopf/Fuß der ersten Seite exakt auf den Satzspiegel ausrichten. Erst bei
% \begin{document} setzen, weil KOMA diese Längen vorher neu berechnet.
\AtBeginDocument{%
  \setplength{firstheadwidth}{\textwidth}%
  \setplength{firstheadhpos}{\dimexpr 1in+\hoffset+\oddsidemargin\relax}%
  \setplength{firstfootwidth}{\textwidth}%
  \setplength{firstfoothpos}{\dimexpr 1in+\hoffset+\oddsidemargin\relax}%
}
\setplength{firstfootvpos}{277mm}

\setkomavar{firsthead}{%
  \begin{minipage}[b]{0.52\textwidth}
    {\NameSchrift\fontsize{18pt}{20pt}\selectfont\color{akzent}%
      \usekomavar{fromname}}%
  \end{minipage}%
  \hfill
  \begin{minipage}[b]{0.44\textwidth}
    \raggedleft\footnotesize\color{grau}%
    \usekomavar{fromaddress}\\[1pt]
    \usekomavar{fromphone}\\[1pt]
    \usekomavar{fromemail}%
  \end{minipage}\\[2.5mm]
  {\color{akzent}\rule{\textwidth}{0.6pt}}%
}

% Erste Seite: Seitenzahl ebenfalls unten links (nur wenn mehrseitig)
\setkomavar{firstfoot}{\footnotesize\color{grau}\seitenzahl\hfill}

% --- Bezugszeichen: nur setzen, wenn es eins gibt (pro Brief) --------------
% z. B. \setkomavar{yourref}{AB-12345}     -> "Ihr Zeichen"
%       \setkomavar{customer}[Aktenzeichen:~]{12345}

% --- Grußformel und Unterschrift linksbündig am Satzspiegel ----------------
\setplength{sigindent}{0mm}
\renewcommand*{\raggedsignature}{\raggedright}
% Abstand Gruß -> Unterschrift. Der Platz für die Unterschrift selbst
% (Bild oder 15 mm Leerraum) kommt aus \unterschriftsbild in absender.lco.
\setplength{sigbeforevskip}{2mm}

\setkomavar{date}{\today}

% --- Anlagen: PDF oder Bild (JPG/PNG) hinten an den Brief anhängen ---------
% Nach \end{letter}, vor \end{document}:
%   \anlage[Mietvertrag]{anlagen/mietvertrag.pdf}   alle Seiten eines PDFs
%   \anlage[Ausweis]{anlagen/ausweis.jpg}            Foto oder Scan
%   \anlage[Rechnung][pages={1,3}]{anlagen/re.pdf}   nur bestimmte Seiten
% Jede Seite bekommt oben rechts den Stempel „Anlage N: <Titel>“.
% Die Anlagenliste im Brief selbst bleibt \encl{...} vor \end{letter}.
\usepackage{pdfpages}
\newcounter{anlage}
\NewDocumentCommand{\anlage}{o O{} m}{%
  \stepcounter{anlage}%
  \def\anlagestempel{Anlage~\theanlage\IfValueT{#1}{: #1}}%
  \IfFileExists{#3}{%
    \includepdf[pages=-, scale=0.92, #2,
      pagecommand={\thispagestyle{empty}},
      picturecommand={\put(\LenToUnit{\paperwidth-20mm},\LenToUnit{\paperheight-12mm}){%
        \makebox[0pt][r]{\footnotesize\sffamily\color{grau}\anlagestempel}}}%
    ]{#3}%
  }{%
    \clearpage\thispagestyle{empty}%
    {\color{red}\bfseries \anlagestempel\ – Datei „\detokenize{#3}“ nicht gefunden}%
    \clearpage
  }%
}
 geladen. Gefunden wird es über
%  _vorlage/latexmkrc (setzt die Suchpfade für TeX-Dateien und Schriften).
%  Schriften: _vorlage/fonts/ · Absenderdaten: <Person>/absender.lco
% ============================================================
\documentclass[
  fontsize=11pt,
  paper=a4,
  parskip=half,
  enlargefirstpage,
  foldmarks=true,
  foldmarks=TBMP,        % Faltmarken oben/Mitte/unten + Lochmarke
  backaddress=true,      % Rücksendeangabe im Fenster
  subject=beforeopening,
  subject=untitled,      % Betreff fett, ohne "Betreff:" (DIN 5008)
]{scrlttr2}

\LoadLetterOption{DIN}   % Anschriftfeld/Geometrie für DIN-lang-Fensterkuvert

% --- Satzspiegel nach DIN 5008: links 25 mm, rechts 20 mm ------------------
% (KOMA-Standard ergab ~31,5 mm links/rechts; Anschrift und Text standen
%  dadurch nicht auf derselben Fluchtlinie.)
\setlength{\oddsidemargin}{\dimexpr 25mm-1in\relax}
\setlength{\evensidemargin}{\oddsidemargin}
\setlength{\textwidth}{\dimexpr\paperwidth-25mm-20mm\relax}
% Anschrift beginnt bei 20 mm (Fensterrand) + 5 mm Einzug = 25 mm,
% also bündig mit Brieftext, Kopf und Betreff.
\setplength{toaddrindent}{5mm}
% Unterer Rand der Folgeseiten: KOMA-Standard ließ den Text bei ca. 239 mm
% enden (fast 6 cm Rand). Jetzt endet er bei 265 mm, die Fußzeile steht bei
% 277 mm. (Seite 1 wird über enlargefirstpage bis kurz vor den Fuß genutzt.)
\setlength{\textheight}{\dimexpr 265mm-(1in+\voffset+\topmargin+\headheight+\headsep)\relax}

% --- Grafiken (eingescannte Unterschrift) — vor absender.lco laden! ------
\usepackage{graphicx}

% --- Absenderdaten: nächstgelegene absender.lco -----------------------------
% Gesucht wird vom Ordner des Briefes aus nach oben. So gilt Max/absender.lco
% für alles unter Max/, Erika/absender.lco für alles unter Erika/.
% Ein Brief kann eine eigene absender.lco im eigenen Ordner haben (gewinnt).
% Nur ein Absender? Dann genügt eine absender.lco im Projektordner.
\directlua{
  local d = lfs.currentdir():gsub(string.char(92), "/")  % Windows: \ -> /
  local found = ""
  while d and d ~= "" do
    if lfs.isfile(d .. "/absender.lco") then found = d break end
    local parent = d:match("^(.*)/[^/]+$")
    if not parent or parent == "" or parent == d then break end
    d = parent
  end
  token.set_macro("absenderdir", found)
}
\ifx\absenderdir\empty
  \errmessage{Keine absender.lco gefunden: Lege eine absender.lco in den
    Ordner des Briefes oder einen Ordner darüber (Vorlage:
    _vorlage/absender-vorlage.lco)}
\else
  \input{\absenderdir/absender.lco}
\fi

% --- Schrift: Source Sans 3 aus ./fonts/ ---------------------------------
\usepackage{fontspec}
% Schriften werden per Dateiname gefunden (Suchpfad aus latexmkrc),
% darum hier KEIN Path=fonts/ – so klappt es aus jedem Unterordner.
\setmainfont{SourceSans3}[
  Extension      = .otf,
  UprightFont    = *-Regular,
  ItalicFont     = *-It,
  BoldFont       = *-Semibold,       % Semibold als \textbf: ruhiger als Bold
  BoldItalicFont = *-SemiboldIt,
  Scale = 1.08,
]
\setsansfont{SourceSans3}[
  Extension=.otf,
  UprightFont=*-Regular, ItalicFont=*-It,
  BoldFont=*-Semibold, BoldItalicFont=*-SemiboldIt,
]% auch \sffamily (z. B. Rücksendeangabe) nutzt Source Sans
\newfontfamily\NameSchrift{PlayfairDisplaySC-Black}[Extension=.ttf]

\linespread{1.06}

\usepackage[ngerman]{babel}
\usepackage{microtype}

% Kompakte Aufzählungen (z. B. Fragenkataloge)
\usepackage{enumitem}
\setlist{itemsep=0pt, parsep=0pt, topsep=2pt, partopsep=0pt, leftmargin=1.6em}

% --- Farben ---------------------------------------------------------------
\usepackage{xcolor}
\definecolor{akzent}{RGB}{0, 82, 100}    % dunkles Petrol
\definecolor{grau}{RGB}{88, 92, 96}      % Text im Kopf/Fuß

% --- PDF-Metadaten + anklickbare Links (ohne farbige Rahmen) ---------------
% Titel setzt jeder Brief selbst: \hypersetup{pdftitle={...}}
\usepackage[hidelinks]{hyperref}
\hypersetup{pdfauthor={\absendername}, pdflang={de-DE}}

% --- Seitenzahl "1/2" — nur bei mehrseitigen Briefen (DIN 5008) ------------
% scrlttr2 schreibt für jeden Brief das Label "<Briefnr>.lastpage".
% \thisletter liefert die Nummer des aktuellen Briefes, damit funktioniert
% die Zählung auch, wenn mehrere Briefe in einer Datei stehen.
% refcount liest die Seitenzahl als Zahl aus; ist der Brief einseitig,
% bleibt die Fußzeile leer.
\usepackage{refcount}
\newcommand*{\seitenzahl}{%
  \ifnum\getpagerefnumber{\thisletter.lastpage}>1
    \thepage/\pageref{\thisletter.lastpage}%
  \fi}

% Folgeseiten: Seitenzahl unten LINKS (KOMA-Mechanismus scrlayer-scrpage)
\usepackage[automark]{scrlayer-scrpage}
\clearpairofpagestyles
\ifoot*{\footnotesize\color{grau}\seitenzahl}
\pagestyle{scrheadings}

% ==========================================================================
%  BRIEFKOPF: Name links, Kontaktblock rechts, feine Akzentlinie darunter.
%  Gespeist aus den Variablen in absender.lco – hier nichts eintragen.
% ==========================================================================
% Kopf/Fuß der ersten Seite exakt auf den Satzspiegel ausrichten. Erst bei
% \begin{document} setzen, weil KOMA diese Längen vorher neu berechnet.
\AtBeginDocument{%
  \setplength{firstheadwidth}{\textwidth}%
  \setplength{firstheadhpos}{\dimexpr 1in+\hoffset+\oddsidemargin\relax}%
  \setplength{firstfootwidth}{\textwidth}%
  \setplength{firstfoothpos}{\dimexpr 1in+\hoffset+\oddsidemargin\relax}%
}
\setplength{firstfootvpos}{277mm}

\setkomavar{firsthead}{%
  \begin{minipage}[b]{0.52\textwidth}
    {\NameSchrift\fontsize{18pt}{20pt}\selectfont\color{akzent}%
      \usekomavar{fromname}}%
  \end{minipage}%
  \hfill
  \begin{minipage}[b]{0.44\textwidth}
    \raggedleft\footnotesize\color{grau}%
    \usekomavar{fromaddress}\\[1pt]
    \usekomavar{fromphone}\\[1pt]
    \usekomavar{fromemail}%
  \end{minipage}\\[2.5mm]
  {\color{akzent}\rule{\textwidth}{0.6pt}}%
}

% Erste Seite: Seitenzahl ebenfalls unten links (nur wenn mehrseitig)
\setkomavar{firstfoot}{\footnotesize\color{grau}\seitenzahl\hfill}

% --- Bezugszeichen: nur setzen, wenn es eins gibt (pro Brief) --------------
% z. B. \setkomavar{yourref}{AB-12345}     -> "Ihr Zeichen"
%       \setkomavar{customer}[Aktenzeichen:~]{12345}

% --- Grußformel und Unterschrift linksbündig am Satzspiegel ----------------
\setplength{sigindent}{0mm}
\renewcommand*{\raggedsignature}{\raggedright}
% Abstand Gruß -> Unterschrift. Der Platz für die Unterschrift selbst
% (Bild oder 15 mm Leerraum) kommt aus \unterschriftsbild in absender.lco.
\setplength{sigbeforevskip}{2mm}

\setkomavar{date}{\today}

% --- Anlagen: PDF oder Bild (JPG/PNG) hinten an den Brief anhängen ---------
% Nach \end{letter}, vor \end{document}:
%   \anlage[Mietvertrag]{anlagen/mietvertrag.pdf}   alle Seiten eines PDFs
%   \anlage[Ausweis]{anlagen/ausweis.jpg}            Foto oder Scan
%   \anlage[Rechnung][pages={1,3}]{anlagen/re.pdf}   nur bestimmte Seiten
% Jede Seite bekommt oben rechts den Stempel „Anlage N: <Titel>“.
% Die Anlagenliste im Brief selbst bleibt \encl{...} vor \end{letter}.
\usepackage{pdfpages}
\newcounter{anlage}
\NewDocumentCommand{\anlage}{o O{} m}{%
  \stepcounter{anlage}%
  \def\anlagestempel{Anlage~\theanlage\IfValueT{#1}{: #1}}%
  \IfFileExists{#3}{%
    \includepdf[pages=-, scale=0.92, #2,
      pagecommand={\thispagestyle{empty}},
      picturecommand={\put(\LenToUnit{\paperwidth-20mm},\LenToUnit{\paperheight-12mm}){%
        \makebox[0pt][r]{\footnotesize\sffamily\color{grau}\anlagestempel}}}%
    ]{#3}%
  }{%
    \clearpage\thispagestyle{empty}%
    {\color{red}\bfseries \anlagestempel\ – Datei „\detokenize{#3}“ nicht gefunden}%
    \clearpage
  }%
}
